% ============================================================================
% CAPITULO 3: METODOLOGIA Y DESARROLLO
% ============================================================================

\section{Metodolog\'ia y Desarrollo}
\label{sec:metodologia}

% TODO: Redactar contenido usando skill tfg-cap3-metodologia
% Este es el capitulo mas largo (~20-25 paginas)
% Documenta el sistema ACTUAL antes de cualquier cambio

% ----------------------------------------------------------------------------
% 3.1 Metodologia de Desarrollo
% ----------------------------------------------------------------------------
\subsection{Metodolog\'ia de Desarrollo}
\label{subsec:metodologia-desarrollo}

% Enfoque iterativo incremental (9 fases, 43 dias)
% Herramientas: Python 3.12, Git, Docker, CI/CD
% Proceso: disenar -> implementar -> evaluar -> iterar
% Gestion del proyecto: commits atomicos, documentacion continua (CLAUDE.md)

% ----------------------------------------------------------------------------
% 3.2 Arquitectura del Sistema
% ----------------------------------------------------------------------------
\subsection{Arquitectura del Sistema}
\label{subsec:arquitectura-sistema}

% Diagrama de arquitectura general (web + telegram)
% Stack tecnologico completo
% Componentes principales y sus responsabilidades
% Decisiones de diseno justificadas

% --- 3.2.1 Arquitectura Web ---
\subsubsection{Arquitectura Web (v3.3)}
\label{subsubsec:arquitectura-web}

% FastAPI + WebSocket + frontend HTML/CSS/JS
% Diagrama de flujo: usuario -> WebSocket -> backend -> RAG -> Ollama
% Colores corporativos DNI (#5B7FDB)

% --- 3.2.2 Arquitectura Telegram ---
\subsubsection{Arquitectura Telegram (v4.1.1)}
\label{subsubsec:arquitectura-telegram}

% python-telegram-bot v20+ con polling
% Service layer (4 servicios async)
% PostgreSQL + Alembic migrations
% Diagrama: usuario -> Telegram API -> handlers -> services -> DB + RAG

% ----------------------------------------------------------------------------
% 3.3 Pipeline RAG
% ----------------------------------------------------------------------------
\subsection{Pipeline RAG}
\label{subsec:pipeline-rag}

% --- 3.3.1 Ingesta y Chunking ---
\subsubsection{Ingesta y Chunking de Documentos}
\label{subsubsec:ingesta-chunking}

% 16 documentos DNI -> 263 chunks (197 FAQ + 66 regulares)
% FAQ-aware chunking: deteccion formato Q:/A:
% Script: 02_create_faq_aware_chunks.py
% Embeddings: paraphrase-multilingual-mpnet-base-v2 (768 dim)

% --- 3.3.2 Busqueda Hibrida ---
\subsubsection{B\'usqueda H\'ibrida}
\label{subsubsec:busqueda-hibrida-impl}

% BM25 (50%) + Semantic (50%)
% ChromaDB como vector store
% top_k=10, similarity_threshold=0.30
% ConfigurableRAGEngine (rag_engine.py)

% --- 3.3.3 Generacion y Validacion ---
\subsubsection{Generaci\'on y Validaci\'on Adaptativa}
\label{subsubsec:generacion-validacion}

% EnhancedRAGEngineNew
% question_id=0 forzado (RAG avanzado siempre)
% 10-15 chunks consultados
% Cross-encoder reranking
% Confidence dinamico (6 factores)
% System prompt v3.2 completo

% ----------------------------------------------------------------------------
% 3.4 Context Tracking
% ----------------------------------------------------------------------------
\subsection{Context Tracking Inteligente}
\label{subsec:context-tracking}

% --- 3.4.1 ContextTracker ---
\subsubsection{ContextTracker: Detecci\'on de Proyectos DNI}
\label{subsubsec:context-tracker-impl}

% context_tracker.py (415 lineas)
% 6 proyectos detectables con scoring ponderado
% Ventana deslizante 4 interacciones
% Ponderacion por recencia (user 1.0, assistant 0.15)
% Umbral de confianza: 0.6

% --- 3.4.2 Query Enrichment ---
\subsubsection{Query Enrichment Autom\'atico}
\label{subsubsec:query-enrichment-impl}

% Prefijo contextual vs reformulacion agresiva
% "[Contexto: Desayunos Solidarios] query"
% Extraccion de tema conversacional (6 temas)
% Multi-factor confidence: project(0.7) + topic(0.3)

% --- 3.4.3 ConversationalRAG ---
\subsubsection{ConversationalRAG Engine}
\label{subsubsec:conversational-rag}

% conversational_rag.py (492 lineas)
% Integracion ContextTracker + RAG engine
% Historial de conversacion por sesion
% Fallback mejorado (10 chunks, temperature=0.2)

% ----------------------------------------------------------------------------
% 3.5 Persistencia Cross-Sesion (Telegram)
% ----------------------------------------------------------------------------
\subsection{Persistencia Cross-Sesi\'on}
\label{subsec:persistencia-cross-sesion}

% --- 3.5.1 Esquema de Base de Datos ---
\subsubsection{Esquema de Base de Datos PostgreSQL}
\label{subsubsec:esquema-bd}

% 7 tablas: users, conversations, messages, context_snapshots, feedback,
%           user_consents, analytics_events
% Diagrama ER
% JSONB para metadata flexible
% Indices optimizados
% Alembic migrations

% --- 3.5.2 PersistentContextTracker ---
\subsubsection{PersistentContextTracker}
\label{subsubsec:persistent-context-tracker}

% persistent_context_tracker.py (345 lineas)
% Exponential decay (half-life = 3 dias)
% Merge de contextos: recent + historical
% Triggers de snapshots (cada 5 mensajes, cambio proyecto, high confidence)
% Sistema anti-contaminacion v4.1.1 (3 capas)

% --- 3.5.3 Service Layer ---
\subsubsection{Service Layer}
\label{subsubsec:service-layer}

% 4 servicios async: UserService, ConversationService,
%   MessageService, ContextService
% Patron Repository con SQLAlchemy ORM
% expire_on_commit=False (fix critico)

% ----------------------------------------------------------------------------
% 3.6 Ensemble Multi-Modelo
% ----------------------------------------------------------------------------
\subsection{Sistema Ensemble Multi-Modelo}
\label{subsec:ensemble-impl}

% --- 3.6.1 Arquitectura Ensemble ---
\subsubsection{Arquitectura del Motor Ensemble}
\label{subsubsec:arquitectura-ensemble}

% ensemble_engine.py (287 lineas)
% question_classifier.py (130 lineas)
% 4 estrategias en strategies/

% --- 3.6.2 Estrategias ---
\subsubsection{Implementaci\'on de las 4 Estrategias}
\label{subsubsec:estrategias-impl}

% Voting (92 lineas), Weighted (109), Routing (184), Consensus (142)
% Resultados: Consensus 0.903, Routing 0.895, Weighted 0.889, Voting 0.872
% Decision: gemma2:27b directo en produccion (ensemble disponible)

% ----------------------------------------------------------------------------
% 3.7 Seguridad
% ----------------------------------------------------------------------------
\subsection{Seguridad}
\label{subsec:seguridad}

% --- 3.7.1 Amenazas Identificadas ---
\subsubsection{Amenazas Identificadas}
\label{subsubsec:amenazas}

% Prompt injection, XSS, CORS, input validation
% OWASP Top 10 aplicado a chatbots

% --- 3.7.2 Medidas Implementadas ---
\subsubsection{Medidas Implementadas}
\label{subsubsec:medidas-seguridad}

% SDK de seguridad: InjectionDetector, Sanitizer, RiskScorer, RateLimiter
% Integracion en backend web y Telegram
% GDPR compliance (/delete_my_data)

% ----------------------------------------------------------------------------
% 3.8 Infraestructura y Deployment
% ----------------------------------------------------------------------------
\subsection{Infraestructura y Deployment}
\label{subsec:infraestructura}

% --- 3.8.1 Servidor Ollama UPV ---
\subsubsection{Servidor Ollama UPV}
\label{subsubsec:servidor-ollama}

% ollama.gti-ia.upv.es:443
% gemma2:27b como modelo principal
% Sin costos, privacidad, control total

% --- 3.8.2 Docker y PostgreSQL ---
\subsubsection{Docker y PostgreSQL}
\label{subsubsec:docker-postgres}

% docker-compose.yml: PostgreSQL 16 en puerto 5434
% Alembic para migrations
% Volumes persistentes

% --- 3.8.3 CI/CD ---
\subsubsection{Integraci\'on Continua}
\label{subsubsec:cicd}

% GitHub Actions: 5 workflows
% Testing automatizado
% Linting y type checking
